%
% troposkein.tex -- links haengende Schnur (Kettenlinie), rechts
% rotierendes Troposkein y(x) = a*sn(x/c, k) (Daten aus
% troposkein-werte.dat, RK4, k = 0.55).
%
\documentclass[tikz]{standalone}
\usepackage{amsmath}
\usepackage{times}
\usepackage{txfonts}
\usetikzlibrary{arrows,math,calc}
\input{farben.tex}
\begin{document}
\def\skala{1.02}
\begin{tikzpicture}[>=latex, thick, scale=\skala]

\def\L{2.30}

% ---------- linkes Panel: haengende Schnur ----------
\def\acat{2.6778}      % a in a*(cosh(x/a)-1), Durchhang 1.05 bei x = L

\draw[dashed, black!45] (-\L-0.3,0) -- (\L+0.3,0);
\draw[darkred, line width=1.8pt, domain=-\L:\L, samples=60, smooth]
	plot (\x, {\acat*cosh(\x/\acat)-\acat*cosh(\L/\acat)});

\fill (-\L,0) circle (0.035);
\fill (\L,0) circle (0.035);
\node[anchor=north] at (-\L,-0.05) {$A$};
\node[anchor=north] at (\L,-0.05) {$B$};

% ---------- rechtes Panel: rotierendes Troposkein ----------
\begin{scope}[xshift=6.6cm]

\draw[dashed, black!45] (-\L-0.45,0) -- (\L+0.45,0);

\draw[->, black!55] (\L+0.62,0.28) arc[start angle=70, end angle=-250, radius=0.28];
\node[black!55] at (\L+0.62,0.68) {$\omega$};

\draw[darkred, line width=1.8pt] plot file {troposkein-werte.dat};

\fill (-\L,0) circle (0.035);
\fill (\L,0) circle (0.035);
\node[anchor=north] at (-\L,-0.05) {$A$};
\node[anchor=north] at (\L,-0.05) {$B$};

% Kraeftebilanz an einem Schnurelement bei P: Tangentenwinkel theta,
% Spannung T mit Komponenten T cos(theta) (axial) und T sin(theta)
% (quer), sowie die radial nach aussen wirkende Zentrifugalkraft.
% Die Tangente zeigt nach hinten-oben (Winkel 180-thP), damit T und
% seine Komponenten in die freie Flaeche zwischen P und A zeigen.
\def\Px{-1.188}\def\Py{0.756}
\def\thP{26.6}
\def\rP{180-\thP}
\def\bis{(180+\rP)/2}
\fill (\Px,\Py) circle (0.028);

% Tangente als Referenzlinie
\draw[black!45, thin] (\Px,\Py) -- ++({0.45*cos(\rP)},{0.45*sin(\rP)});
\draw[black!45, thin] (\Px,\Py) -- ++(-0.45,0);

% Winkel theta gegen die Achse
\draw[black!45] ($(\Px,\Py)+({0.28*cos(\rP)},{0.28*sin(\rP)})$)
	arc[start angle=\rP, end angle=180, radius=0.28];
\node[black!45] at ({\Px+0.58*cos(\bis)},{\Py+0.58*sin(\bis)}) {\small $\theta$};

% Spannung T (entlang der Tangente) und ihre Komponenten
\draw[->, black!75, thick] (\Px,\Py) -- ++({1.05*cos(\rP)},{1.05*sin(\rP)})
	node[above=2pt] {\small $T$};
\draw[dashed, black!40] (\Px,\Py) -- ++({1.05*cos(\rP)},0)
	node[midway, below=4pt] {\small $T\cos\theta$};
\draw[dashed, black!40] ({\Px+1.05*cos(\rP)},{\Py}) -- ++(0,{1.05*sin(\rP)})
	node[midway, left=4pt] {\small $T\sin\theta$};

% Zentrifugalkraft, radial nach aussen (in +y Richtung)
\draw[->, darkgreen, thick] (\Px,\Py) -- ++(0,0.85)
	node[above right=-1pt] {\small $\rho\omega^2 y\,ds$};

\end{scope}

\end{tikzpicture}
\end{document}
